% v11 Background and Related Work (L1-L5, I2, REF1-REF3, REF5). Replaces v10 Section II.
\section{Background and Related Work}\label{sec:related}

\subsection{How Often Review Comments Outlive Their Lines}\label{sec:motivation}
To quantify the problem in a public setting, we collected every merged pull request with at least one inline review thread from Requests, Flask, Click, HTTPX and Rich between January 2022 and June 2026, through GitHub's GraphQL interface \cite{githubgraphql}. Of \wpeThreads{} inline threads, \wpeFollowedPct\% were followed by at least one further commit to the pull request, and \wpeOutdatedPct\% (pull-request cluster bootstrap 95\% interval \wpeOutdatedLo--\wpeOutdatedHi\%) were marked outdated; \wpeNullLinePct\% of the outdated threads no longer had a line in the current diff. GitHub marks a thread outdated when the diff hunk it was written on changes \cite{github2018outdating}; outdated threads remain readable, but the reviewer must find where the concern now applies. The sample is professional open source, the snapshot reflects force-pushes and squashes, and resolution does not show that a concern was fixed, so the figures motivate rather than measure the educational problem.

\subsection{Robust Annotation Anchoring}
Robust document annotation predates this system. Brush et al.\ examined annotation positioning after document edits \cite{brush2001robust}; Phelps and Wilensky developed robust intra-document locations \cite{phelps2000robust}. Hypothesis describes quote-based fuzzy anchoring \cite{csillag2013fuzzy}, and the W3C Web Annotation Data Model specifies text-quote and text-position selectors \cite{w3c2017annotation}. For code, recent systems keep annotations attached as programs evolve: Meta-Manager records the provenance of edits in the editor \cite{horvath2024metamanager}, and Magic Markup re-tags document-external markup with a language model \cite{misback2024magic}. These systems target a single developer's working copy and its edit stream; \sys{} receives whole submitted versions without an edit stream and must decide, for each earlier concern, whether its supervisor needs to look again.

\subsection{Code Review and Warning Tracking}
Modern review combines line-specific discussion with human judgement \cite{bacchelli2013expectations,sadowski2018modern}. Review tools differ in how comments survive a revision, and a fixed-line model describes none of them fully. GitHub keeps a comment on its line while the line is unchanged and marks it outdated when the exact anchoring line changes \cite{github2018outdating}. Gerrit ports unresolved comments to a newer patch set by calculating their position in the new file, and attaches them at file level when no position can be calculated; ported comments look like ordinary ones, and resolved comments are not ported \cite{gerritporting}. GitLab can be configured to resolve diff threads automatically when a push makes them outdated \cite{gitlabautoresolve}, the opposite of the human-only resolution adopted here. CodeReviewQA tests whether language models understand review comments by asking them to recognise the change type, localise the change and identify the solution, on 900 curated examples in nine languages \cite{lin2025codereviewqa}. It localises the code a comment asks to change within one version, whereas \sys{} must find where already-reviewed code went in the next. Static-analysis warning trackers retain issue identity across versions and remember earlier review decisions about warnings \cite{avgustinov2015tracking,li2024tracking,spacco2006tracking,sonarqube}; their findings are produced by the analyzer, so an issue that disappears from the analysis can be treated as fixed, whereas a supervisor's concern cannot. ViolationTracker builds precise violation histories with introduced, fixed, deleted and reopened states \cite{yu2023violationtracker}; there too a violation is fixed when the analyzer stops reporting it. Whether a revision touched the reviewed code and whether the concern about that code still holds are different questions; a \sys{} change label answers only the first, so it only requests a check. Table~\ref{tab:related} compares documented behaviours. Apart from GitHub's outdated flag, which the suggestion oracle measures, we did not benchmark these products; our conventional baselines reimplement their documented policies, and product documentation supports the table, not performance claims.

\begin{table*}[t]
\caption{Documented Behaviour of Review Tools and Code-Mapping Methods}\label{tab:related}
\centering\footnotesize
\setlength{\tabcolsep}{3pt}
\begin{tabular}{@{}p{.21\textwidth}p{.12\textwidth}p{.12\textwidth}p{.12\textwidth}p{.12\textwidth}p{.12\textwidth}p{.12\textwidth}@{}}\toprule
Criterion & GitHub \cite{github2018outdating} & Gerrit \cite{gerritporting} & GitLab \cite{gitlabautoresolve} & CodeMapper \cite{hu2026codemapper} & LHDiff-style \cite{asaduzzaman2013lhdiff} & \sys{} \\ \midrule
Follows unchanged code shifted by edits above & yes & calculated position & not documented & yes & yes & yes \\
Follows code edited in place & no: marked outdated & calculated position or file level & no: outdated & yes & yes & yes, if similarity $\ge\theta$ \\
Follows code moved to another file & not documented & not documented & not documented & via file renames & no & yes, exact text \\
Change status shown to reviewer & outdated & none (looks unchanged) & outdated & n/a (mapping only) & n/a & content-match, moved, changed, removed \\
Who resolves & people & people & people, or automatically when outdated (option) & n/a & n/a & supervisor only \\
Dismissal remembered on later versions & not documented & not documented & not documented & n/a & n/a & yes, by finding key \\
Published relocation accuracy & none found & none found & none found & yes & yes & this paper \\ \bottomrule
\end{tabular}
\tablenote{}{Product cells summarize the cited documentation as accessed on 31 October 2026, not measured behaviour. ``Not documented'': no statement found in the vendor's documentation. ``None found'': no published evaluation of the product's relocation accuracy was found in its documentation or in our related-work search. n/a: the method maps code regions and has no review state.}
\end{table*}

\subsection{Source Differencing and Region Mapping}
Line differencing underlies many revision tools \cite{myers1986ond}. Our matcher uses Python's sequence matcher, which extends the Ratcliff--Obershelp approach \cite{ratcliff1988pattern,pythonsequence}. LHDiff combines line content and context to track source lines \cite{asaduzzaman2013lhdiff}; related work tracks changed lines \cite{canfora2007identifying}, surveys program-element matching \cite{kim2006program} and tracks source locations \cite{reiss2008tracking}. ChangeDistiller and GumTree provide syntax-aware differencing \cite{fluri2007change,falleri2014fine}; a refactoring-aware syntax-tree differencer built on RefactoringMiner adds multi-mappings and commit-level diffs, with a benchmark of node mappings from bug-fixing and refactoring commits \cite{alikhanifard2024astdiff}. Origin analysis considers renamed, merged and split entities \cite{godfrey2005using}. History trackers follow named program elements through a commit history with refactoring awareness: CodeShovel builds method histories \cite{grund2021codeshovel}, CodeTracker tracks methods and variables \cite{jodavi2022codetracker} and code blocks \cite{hasan2024blocktracking}, and HistoryFinder revisits the oracles these tools were evaluated on \cite{islam2025historyfinder}. They track syntactic elements, whereas a review anchor is an arbitrary range of lines. ReCite finds Linux-kernel comments whose referenced function was renamed or removed, traces the reference through the history and proposes a repair with a language model for people to accept \cite{sun2026recite}. The closest recent work is CodeMapper, which maps an arbitrary code region from one commit to another, language-independently, by combining several git diff algorithms, movement detection and text search with context-aware similarity selection \cite{hu2026codemapper}. These methods answer \emph{where is this code now}. None defines what a reviewer should be told about the change, or what happens to a concern when the answer is uncertain; that is the part \sys{} adds, and Section~\ref{sec:external} compares the matchers directly.

\subsection{Issue Lifecycles}
Issue trackers, review discussions and warning trackers all carry items across time with states such as open, resolved and closed. The relevant distinction is who may close an item and on what evidence. Analyzer-driven trackers may close an issue the analysis no longer reports; some review tools may close a thread when its lines change \cite{gitlabautoresolve}. \sys{} restricts every closing transition to the supervisor and lets relocation evidence move an issue only towards \lbl{needs re-check}. We do not claim that human-controlled resolution or persistent dismissal is new in isolation; the contribution is binding them to relocation labels with an explicit, checked transition relation (Section~\ref{sec:design}).

\subsection{Educational Feedback}
Feedback research motivates specific, actionable comments \cite{hattie2007power,shute2008focus}, and feedback literacy concerns how learners interpret and use them \cite{carless2018feedback,winstone2017supporting}. Keuning et al.\ found that programming-feedback tools support diagnosing mistakes more than planning corrections \cite{keuning2018systematic}. This literature establishes that revisiting feedback matters; it does not establish that relocating annotations improves students' understanding or response to feedback, and we know of no study that does. That causal link is the open question this paper's technical results make testable, not one they answer.
